This chapter reports measured results. All numbers come from the evaluation split unless
labelled \emph{calibration}; paper values are shown only as anchors and are never mixed into
measured rows.

\section{Experimental Setup}
\begin{itemize}
	\item \textbf{Classes.} Common (COCO scenes), facial (portraits), style (cat photographs).
	\item \textbf{Channel.} SwinJSCC ($C=64$) over AWGN at 5, 10, 15 and 20\,dB SNR.
	\item \textbf{Hiding.} EDICT, 20 steps, strength 0.2, mixing 0.93, $\eta = 0.05$, IPAdapter-plus
	weight 1, $256$\,px.
	\item \textbf{Sample sizes.} The baseline and the RACS comparison were measured on 2 images per
	class at all four SNRs (24 cases per system); the transmitter--receiver invariant and the
	exact-recovery layer were measured on all 66 evaluation images.
\end{itemize}

\section{Baseline Recovery Quality}
\begin{table}[H]
	\centering
	\caption{Baseline recovery PSNR (dB) by class and SNR.}
	\label{tab:basepsnr}
	\begin{tabular}{|l|c|c|c|c|}
		\hline
		\textbf{Class} & \textbf{5 dB} & \textbf{10 dB} & \textbf{15 dB} & \textbf{20 dB} \\ \hline
		Common & 18.30 & 20.72 & 21.53 & 21.45 \\ \hline
		Facial & 22.56 & 24.53 & 24.89 & 25.10 \\ \hline
		Style & 19.93 & 22.68 & 23.28 & 23.43 \\ \hline
	\end{tabular}
\end{table}

Quality rises with SNR and flattens above 15\,dB, where the channel stops being the bottleneck and
the VAE and diffusion round trip limit quality. The perceptual metrics follow the same pattern:
from 5 to 20\,dB, SSIM improves from 0.42--0.60 to 0.58--0.72 and LPIPS falls from 0.23--0.46 to
0.14--0.24 across the three classes. Faces score highest, helped by CodeFormer.
Against the paper's Table~II anchors, the facial class at 5\,dB reaches 22.56\,dB (paper 22.73)
and at 20\,dB 25.10\,dB with LPIPS 0.140 (paper 26.29\,dB, LPIPS 0.213); the common class at
5\,dB reaches 18.30\,dB (paper 19.71); the style class at 5\,dB has LPIPS 0.345 (paper 0.385).
Absolute PSNR is about 1--3\,dB below the paper, consistent with the reduced resolution, step
count and strength forced by the hardware (Section~\ref{sec:deviations}).

\section{Baseline Security and Steganalysis}
\begin{table}[H]
	\centering
	\caption{Recovery PSNR under the three receiver conditions (10\,dB).}
	\label{tab:security}
	\begin{tabular}{|l|c|c|}
		\hline
		\textbf{Condition} & \textbf{This work} & \textbf{Paper (approx.)} \\ \hline
		Correct token & 22.37 dB & 25.37 dB \\ \hline
		Wrong token & 13.96 dB & 18.91 dB \\ \hline
		No token & 14.02 dB & -- \\ \hline
		\textbf{Correct $-$ wrong gap} & \textbf{8.41 dB} & 6.46 dB \\ \hline
		XuNet detection accuracy & 0.500 & 0.509 \\ \hline
	\end{tabular}
\end{table}

A wrong or missing token reduces recovery to about 14\,dB, an unrecognisable image. The gap
between correct and wrong token is larger than in the paper, so a wrong key reveals less. XuNet is
at chance, so this detector cannot distinguish stego images from clean ones.

\section{RACS versus the Baseline}
\subsection{Protocol invariant}
In separate transmitter and receiver processes, RACS selected the same reference for
\textbf{66 of 66} evaluation images, with identical file hashes. The FAISS index and the
in-process exact search returned identical top-50 rows and scores.

\subsection{Efficiency}
Obtaining the reference is the stage RACS replaces, and it is where the difference is largest
(Figure~\ref{fig:eff}). Diffusion generation takes 3.92\,s and 2.61\,GB of GPU memory per
reference; retrieval takes 0.048\,s (0.51\,s including loading the image from disk) and needs no
GPU at all, about $82\times$ faster. Because both ends must obtain the reference for every message,
the saving applies twice per message.

\begin{figure}[H]
	\centering
	\begin{subfigure}{0.48\textwidth}
		\centering
		\begin{tikzpicture}
			\begin{axis}[vizbar, xbar, bar width=12pt, width=\textwidth, height=3.8cm,
				symbolic y coords={RACS,Baseline}, ytick=data, enlarge y limits=0.5,
				xmin=0, xmax=5, xlabel={Seconds per reference},
				nodes near coords, nodes near coords style={font=\scriptsize, /pgf/number format/fixed, /pgf/number format/precision=3}]
				\addplot[fill=vizA, draw=none] coordinates {(3.92,Baseline) (0.048,RACS)};
			\end{axis}
		\end{tikzpicture}
		\caption{Reference latency.}
	\end{subfigure}\hfill
	\begin{subfigure}{0.48\textwidth}
		\centering
		\begin{tikzpicture}
			\begin{axis}[vizbar, xbar, bar width=12pt, width=\textwidth, height=3.8cm,
				symbolic y coords={RACS,Baseline}, ytick=data, enlarge y limits=0.5,
				xmin=0, xmax=3.4, xlabel={Peak GPU memory (GB)},
				nodes near coords, nodes near coords style={font=\scriptsize}]
				\addplot[fill=vizA, draw=none] coordinates {(2.61,Baseline) (0,RACS)};
			\end{axis}
		\end{tikzpicture}
		\caption{Reference GPU memory.}
	\end{subfigure}
	\caption{Cost of obtaining the reference image: retrieval is about $82\times$ faster and needs no GPU.}
	\label{fig:eff}
\end{figure}

\subsection{Quality and security}
Both systems were measured on the same cases, tokens and code path. Figure~\ref{fig:psnrplot}
shows recovery PSNR for every class and SNR: the RACS curve (solid) lies on or slightly above the
baseline curve (dashed) in 11 of the 12 cells, with the largest gain for faces (up to $+0.64$\,dB).

\begin{figure}[H]
	\centering
	\begin{tikzpicture}
		\begin{axis}[
			width=0.85\textwidth, height=6.5cm,
			xlabel={SNR (dB)}, ylabel={Recovery PSNR (dB)},
			xtick={5,10,15,20}, ymin=17, ymax=26.5,
			legend style={font=\scriptsize, at={(0.98,0.03)}, anchor=south east, legend columns=2},
			viz]
			\addplot[vizC, dashed, mark=o] coordinates {(5,18.30) (10,20.72) (15,21.53) (20,21.45)};
			\addlegendentry{Common, baseline}
			\addplot[vizC, thick, mark=*] coordinates {(5,18.45) (10,20.86) (15,21.57) (20,21.55)};
			\addlegendentry{Common, RACS}
			\addplot[vizA, dashed, mark=square] coordinates {(5,22.56) (10,24.53) (15,24.89) (20,25.10)};
			\addlegendentry{Facial, baseline}
			\addplot[vizA, thick, mark=square*] coordinates {(5,23.19) (10,25.07) (15,25.53) (20,25.60)};
			\addlegendentry{Facial, RACS}
			\addplot[vizB, dashed, mark=triangle] coordinates {(5,19.93) (10,22.68) (15,23.28) (20,23.43)};
			\addlegendentry{Style, baseline}
			\addplot[vizB, thick, mark=triangle*] coordinates {(5,19.88) (10,22.83) (15,23.41) (20,23.45)};
			\addlegendentry{Style, RACS}
		\end{axis}
	\end{tikzpicture}
	\caption{Recovery PSNR versus SNR for both systems.}
	\label{fig:psnrplot}
\end{figure}

\begin{table}[H]
	\centering
	\caption{Aggregate comparison (24 paired cases).}
	\label{tab:racsagg}
	\begin{tabular}{|l|c|c|}
		\hline
		\textbf{Metric} & \textbf{Baseline} & \textbf{RACS} \\ \hline
		PSNR (dB) $\uparrow$ & 22.37 & \textbf{22.62} \\ \hline
		SSIM $\uparrow$ & 0.616 & \textbf{0.627} \\ \hline
		LPIPS $\downarrow$ & 0.238 & \textbf{0.227} \\ \hline
		Correct-token PSNR, 10 dB & 22.64 & 22.92 \\ \hline
		Wrong-token PSNR, 10 dB & 13.97 & 13.98 \\ \hline
		No-token PSNR, 10 dB & 14.09 & 14.08 \\ \hline
		Wrong token and wrong reference, 10 dB & -- & 13.98 \\ \hline
		\textbf{Security gap (correct $-$ wrong), 10 dB} & 8.67 dB & \textbf{8.94 dB} \\ \hline
	\end{tabular}
\end{table}

The paired mean PSNR difference is $+0.25$\,dB in favour of RACS, with SSIM and LPIPS also
slightly better and the security gap preserved. A plausible explanation is that a retrieved
photograph is a natural image whose statistics suit the IPAdapter conditioning better than a
generated $256$\,px image does. With 2 images per class this difference is indicative rather than
statistically established; the important finding is that replacing generation with retrieval
does \emph{not} degrade quality or security while removing the GPU cost of the reference.


\subsection{Hypotheses}
\begin{table}[H]
	\centering
	\caption{Outcome of the four hypotheses.}
	\label{tab:hyp}
	\begin{tabularx}{\textwidth}{|l|X|c|}
		\hline
		\textbf{Hypothesis} & \textbf{Evidence} & \textbf{Outcome} \\ \hline
		H1 Efficiency & 82$\times$ faster reference, no GPU memory & Supported \\ \hline
		H2 Quality & PSNR, SSIM, LPIPS equal or better & Supported \\ \hline
		H3 Security & Gap 8.94 vs 8.67 dB & Supported \\ \hline
		H4 Robustness & Same trend across 5--20 dB, all classes & Supported \\ \hline
	\end{tabularx}
\end{table}

\section{RACS Exact-Recovery Layer}
Table~\ref{tab:c2} compares four variants on all 66 evaluation secrets. Variant A writes the
original file into two LSBs; B adds lossless compaction; C replaces LSB writing with
texture-adaptive $\pm1$ embedding; D (proposed) adds the dithered carrier.

\begin{table}[H]
	\centering
	\caption{Exact-recovery layer, 66 secrets (each stego written to and read back from PNG).}
	\label{tab:c2}
	\resizebox{\textwidth}{!}{%
	\begin{tabular}{|l|c|c|c|c|c|c|c|}
		\hline
		\textbf{Variant} & \textbf{Exact} & \textbf{Wrong key} & \textbf{Payload} & \textbf{Bits per} & \textbf{PSNR} & \textbf{Chi-square} & \textbf{SPA} \\
		& & \textbf{rejected} & \textbf{(KB)} & \textbf{value} & \textbf{(dB)} & \textbf{AUC} & \textbf{AUC} \\ \hline
		Generative only & 0\,\% & -- & -- & -- & -- & -- & -- \\ \hline
		A: naive 2-LSB & 100\,\% & 100\,\% & 386 & 2.00 & 44.1 & 1.00 & 1.00 \\ \hline
		B: + compression & 100\,\% & 100\,\% & 261 & 2.00 & 44.1 & 1.00 & 1.00 \\ \hline
		C: + textured $\pm1$ & 100\,\% & 100\,\% & 261 & 0.25 & 57.2 & \textbf{0.50} & 1.00 \\ \hline
		\textbf{D: + dither (RACS)} & \textbf{100\,\%} & \textbf{100\,\%} & \textbf{261} & \textbf{0.25} & \textbf{57.2} & \textbf{0.50} & \textbf{0.83} \\ \hline
	\end{tabular}}
\end{table}

\begin{figure}[H]
	\centering
	\hyphenpenalty=10000\exhyphenpenalty=10000
	\begin{tikzpicture}
		\begin{axis}[vizbar, ybar, bar width=11pt, width=0.82\textwidth, height=5.8cm,
			symbolic x coords={A,B,C,D}, xtick=data,
			xticklabels={A: naive 2-LSB, B: + compression, C: + textured $\pm1$, D: + dither (RACS)},
			x tick label style={font=\scriptsize, text width=2.3cm, align=center},
			ymin=0, ymax=1.12, ylabel={Detector AUC}, enlarge x limits=0.18,
			legend style={at={(0.5,1.02)}, anchor=south, legend columns=2}]
			\addplot[fill=vizA, draw=none] coordinates {(A,1.00) (B,1.00) (C,0.50) (D,0.50)};
			\addplot[fill=vizB, draw=none] coordinates {(A,1.00) (B,1.00) (C,1.00) (D,0.83)};
			\legend{Chi-square, SPA}
			\draw[dashed, vizInk, line width=0.6pt] (rel axis cs:0,0.4464) -- (rel axis cs:1,0.4464)
				node[pos=0.5, above, font=\scriptsize, text=vizInk, fill=white, inner sep=1pt] {chance};
		\end{axis}
	\end{tikzpicture}
	\caption{How well two classical detectors separate stego from clean carriers (0.5 = guessing, 1.0 = always detected). Textured $\pm1$ embedding defeats chi-square; the dithered carrier also weakens SPA.}
	\label{fig:c2auc}
\end{figure}

\begin{itemize}
	\item \textbf{Exactness.} All 66 secrets are recovered pixel-identical by every variant,
	against 22.37\,dB approximate recovery from the generative pipeline alone.
	\item \textbf{Authentication.} All 66 wrong keys are rejected by the MAC; no plausible wrong
	image is ever produced.
	\item \textbf{Compaction} reduces the payload by 32\,\% (386 to 261\,KB), which makes a low
	embedding rate affordable.
	\item \textbf{Chi-square is defeated} by adaptive $\pm1$ embedding (AUC 1.00 to 0.50).
	\item \textbf{SPA is weakened but not defeated.} The estimated embedding rate falls from 0.40
	to 0.09, and the AUC falls from 1.00 to 0.83 only once the carrier carries realistic sensor
	noise; without the dither SPA remains at 1.00 even at 0.25 bits per value. The cover source
	matters as much as the embedding.
	\item \textbf{Distortion} against the payload-free carrier is 57.2\,dB PSNR, compared with
	44.1\,dB for naive embedding.
	\item \textbf{Cost} is about 2.5\,s to embed and 2.2\,s to extract on the CPU.
\end{itemize}

\section{Prototype Performance}
Figure~\ref{fig:app} breaks down one exchange for a $902\times1600$ photograph in the default
hybrid mode. Sending is dominated by the two model stages that give the stego its natural look
(writing the public key and the EDICT pass); the RACS reference costs half a second. Receiving
needs no AI model at all and returns a pixel-identical secret in one to two seconds.

\begin{figure}[H]
	\centering
	\begin{tikzpicture}
		\begin{axis}[vizbar, xbar, bar width=10pt, width=0.8\textwidth, height=6.0cm,
			symbolic y coords={rx,emb,edict,racs,kpub}, ytick={rx,emb,edict,racs,kpub},
			yticklabels={Receive (no model),Encrypt + embed,EDICT stego look,RACS reference,Public key (CPU)},
			y tick label style={font=\small},
			xmin=0, xmax=27, xlabel={Seconds}, enlarge y limits=0.15,
			point meta=explicit symbolic,
			nodes near coords, nodes near coords align={horizontal},
			nodes near coords style={font=\scriptsize, text=black, anchor=west, xshift=2pt}]
			\addplot[fill=vizA, draw=none] coordinates {
				(20,kpub) [$\sim$20 s]
				(0.51,racs) [0.51 s]
				(21,edict) [$\sim$21 s]
				(1.6,emb) [0.7--2.5 s]
				(1.5,rx) [1--2 s]};
		\end{axis}
	\end{tikzpicture}
	\caption{Prototype timings for one photograph (hybrid mode, laptop). Bars show typical values; labels give the measured ranges.}
	\label{fig:app}
\end{figure}

\section{Discussion}
\subsection{Why retrieval does not cost quality}
One might expect a generated reference, built specifically from the secret's own description and
structure, to beat a reference picked from a fixed collection. The measurements say otherwise.
The reference enters the hiding process only through the IPAdapter image embedding at bounded
EDICT strength, so what matters is that it is semantically close and natural-looking, not that it
is pixel-aligned with the secret. A real photograph satisfies this at least as well as a
$256$\,px generated image, and the structural rerank keeps the layout roughly compatible. The
small, consistent gain for faces suggests that real portraits from the corpus suit the IPAdapter
conditioning slightly better than synthetic ones.

\subsection{Security of the token}
The private token does three jobs in RACS: it picks the reference, it seeds the noise and mask in
EDICT, and it derives the keys of the exact layer. An attacker who knows all public information can
rebuild the top-10 candidate pool and so guess the right reference with probability $1/10$; but
with the wrong token the noise and mask are still wrong, and the measured recovery with a wrong
token \emph{and} a wrong reference is the same 13.98\,dB as with a wrong token alone. For the exact
layer, every guess of the 8-digit token costs 200{,}000 PBKDF2 iterations, so trying all $10^8$
tokens costs $2\times10^{13}$ HMAC evaluations, and every wrong guess is rejected by the seal
rather than producing a plausible image that could leak information. A longer key would raise
this cost further at no extra cost to legitimate users.

\subsection{When RACS may not fit}
Retrieval can only return what the corpus contains. For a secret far from every corpus image, the
top-10 pool will be semantically weaker and the stego may look less natural. The 1{,}000-image
corpus was sufficient for the three evaluation classes here; a deployment would need a larger and
more diverse collection, or a retrieve-or-generate fallback that uses diffusion only when no corpus
image is close enough. The exact layer, for its part, requires a lossless channel: it suits file
transfer, not photo sharing that recompresses images.

\section{Deviations from the Paper}\label{sec:deviations}
\begin{table}[H]
	\centering
	\caption{Documented differences from the paper and their reason.}
	\label{tab:dev}
	\begin{tabularx}{\textwidth}{|X|X|X|}
		\hline
		\textbf{Paper} & \textbf{This project} & \textbf{Reason} \\ \hline
		512 px & 256 px & Stable Diffusion + ControlNet exceeds 6 GB at 512 px \\ \hline
		EDICT 50 steps, strength 0.5 & 20 steps, strength 0.2 & Run time and memory; a configuration value \\ \hline
		SwinJSCC $C=48$, CelebA-HQ & $C=64$, high-resolution release & Nearest public checkpoint \\ \hline
		MDJCM digital branch & Not reproduced & No public checkpoint; no digital result claimed \\ \hline
		Unnamed LLM & Qwen2.5-0.5B-Instruct & Local, deterministic, fits memory \\ \hline
		OpenPose on all faces & MediaPipe on 3 of 22 & OpenPose finds no body on tight crops \\ \hline
		CodeFormer on GPU & On CPU & GPU memory \\ \hline
		XuNet details & Public XuNet design, image-grouped patches & Not fully specified; small test set \\ \hline
	\end{tabularx}
\end{table}

\section{Limitations}
\begin{itemize}
	\item Quality comparisons use 2 images per class (24 cases per system); the full 22-per-class
	run is a configuration change requiring several hours.
	\item Detectability of the exact layer is measured against classical detectors only; a trained
	CNN steganalyser on this cover source remains to be evaluated.
	\item The exact layer is fragile by design: recompression destroys the payload (detected, not
	silently accepted), so the stego must be sent as a file.
	\item The 8-digit token is short; key stretching slows brute force, but a longer key is
	advisable in practice.
	\item One style evaluation image is a human portrait rather than a cat photograph.
\end{itemize}
